\documentclass{article}\begin{document}At each positive length and base object, changing the terminal Cat-valued payload along an ordinary natural transformation T and then evaluating equals evaluating and then applying the component of layerMap T. This is EvalPayload.evalSquare, an equality of actual functors on the native weighted span fibre, not merely agreement on objects or a comparison transported from an opaque equivalence. On arbitrary chain-total arrows, rawMap retains the endpoint correction of payload naturality; alignedMap uses the all-arrow square comparison and the total-map correction. totalSquare and nativeSquare express equality before restricting to vertical arrows. strictSquare restricts the resulting square to strict fibres, and onFibreSquare handles the native object presentation. embedSquare and valueSquare identify the strict-fibre decoder comparison with the actual component of T, using the Grothendieck inclusion-map comparison and both inverse functor laws for embedding and decoding. The result uses the existing split endpoints and the indicated grade and payload universes. No identity hypothesis on native comparison maps is used. The remaining mixed equation compares the two pastes obtained by traversing \#189's base evaluation square before or after \#190's two chosen base/payload faces, conjugated by these payload comparison equalities. Equality of underlying total functors supplies a common underlying path but does not itself identify those factorizations through selected cartesian lifts. A prospective proof must compare their postcomposites with the same target lift and use uniqueness, then transport through the decoder. That equation, simultaneous coherence, length zero, coproduct assembly, flattening, and strict algebra are not proved here. Manuscript 458ff36a, lines 929--953.\end{document}
